\scanpage{88}
\mspage{43}
à la condition suivante

(123) (Simp) Pour $\forall\ G\in\unsure{\mathcal M|\mathfrak g}$, l'ens. ordonné des quotients \inserted{$G'$} de $G$ est \deleted{isomorphe à un \unsure{simplexe}}\inserted{un « simplexe », i.e. isomorphe} \unsure{à} un ens. des parties d'un ens. fini.

Si $A_G$ désigne l'ens. des quotients de $G$ qui sont $\ne G$ et qui sont \deleted{minimaux}\inserted{maximaux} parmi ceux-ci, alors l'ens. \inserted{ordonné} des quotients de $G$ est donc isomorphe à l'ens. des parties de $A_G$. \deleted{Si} On désigne par $G/C$ le quotient de $G$ associé à $A_G\setminus C$ (d'autant plus petit que $C$ est plus grand).

\margin{\uline{NB} Il faudrait montrer comment cette axiomatisation s'applique à \unsure{$\forall$|la} situation de multiplicités, \deleted{topologiques ou \unsure{étales}}\inserted{\unsure{analytiques relatives|algébriques relatives}} avec un diviseur à croisements normaux\ldots (à cela près que $S\Pi M_G$ est une extension de $\Pi M_G$ non pas par un $\mathbf Z^C$ tordu, mais par un quotient d'un tel\ldots}

On suppose donné pour $\forall\ G\in\operatorname{Ob}\mathcal G$ deux groupoïdes $S\Pi M_G$ et $\Pi M_G$, et un morphisme
\[
\tag{124}
\diagram{\begin{tikzcd}[ampersand replacement=\&]
S\Pi M_G \arrow[r,dash,"\varphi_G"] \& \Pi M_G.
\end{tikzcd}}
\]
On suppose que $\varphi_G$ fait de $S\Pi M_G$ une extension de $\Pi M_G$ par $\mathbf Z^{A(G)}$ (en fait, c'est une structure supplémentaire sur $S\Pi M_G$). On suppose que $G\mapsto\Pi M_G$ est un foncteur en $G$ \uuline{pour les iso} (pas pour les morphismes quelc.).
\margin{87}
